\documentclass[a4paper,14pt]{extarticle}
\usepackage{styles}
\usepackage{svg}

\begin{document}

\litSubsubsection{2.2.1 CUDA Architecture Overview and Fitness Kernel Design}

This section covers CUDA architecture and kernel design concepts needed to understand GPU-accelerated genetic algorithms.

\subsubsection*{CUDA Architectural Overview}

In 2006, NVIDIA launched the first CUDA-capable GPU (GeForce 8800 GTX) along with the CUDA C compiler~\cite{sanders2010cuda}, enabling general-purpose GPU programming. The CUDA Toolkit supports development on systems ranging from embedded devices to HPC supercomputers~\cite{nvidia2024toolkit}.

\begin{figure}[H]
    \centering
    \includegraphics[width=0.8\textwidth]{diagrams/sec2/cpu-gpu.pdf}
    \caption{CPU/GPU Architecture Comparison}
    \label{fig:cpu-gpu-arch}
\end{figure}

\textbf{Hardware Organisation.} A CUDA-capable GPU contains an array of Streaming Multiprocessors (SMs), each with multiple CUDA cores. ``A GPU is built around an array of Streaming Multiprocessors (SMs),'' with programs partitioned into blocks of threads that execute independently across these multiprocessors~\cite{nvidia2025cudaguide}. Each SM contains a register file, shared memory/L1 cache, and functional units~\cite{nvidia2025cudaguide}. Thread blocks are scheduled among available SMs and may execute in any order, enabling automatic load balancing~\cite{hwu2022pmpp}. Figure~\ref{fig:cpu-gpu-arch} illustrates this CPU/GPU architectural difference.

\subsubsection*{What Is a Kernel?}

A CUDA kernel is a function executed on the GPU: ``when called, is executed $N$ times in parallel by $N$ different CUDA threads, as opposed to only once like regular functions''~\cite{nvidia2025cudacpp}. Each thread executes the same code on different data, following the Single-Instruction Multiple-Threads (SIMT) model~\cite{nvidia2025cudaguide}.

\textbf{Syntax and Launch.} In CUDA C/C++, kernels use the \texttt{\_\_global\_\_} qualifier:

{\small
\begin{verbatim}
__global__ void fitness_kernel(float* population, float* fitness, int n) {
    int idx = blockIdx.x * blockDim.x + threadIdx.x;
    if (idx < n)
        fitness[idx] = evaluate_individual(population[idx]);
}
fitness_kernel<<<num_blocks, 256>>>(d_population, d_fitness, pop_size);
\end{verbatim}}

In Numba Python, the equivalent uses \texttt{@cuda.jit} with square brackets for launch configuration:

{\small
\begin{verbatim}
@cuda.jit
def fitness_kernel(population, fitness, n):
    idx = cuda.blockIdx.x * cuda.blockDim.x + cuda.threadIdx.x
    if idx < n:
        fitness[idx] = evaluate_individual(population[idx])

fitness_kernel[num_blocks, 256](d_population, d_fitness, pop_size)
\end{verbatim}}

This work uses Numba for all GPU implementations~\cite{numba2025cuda}.

\subsubsection*{Thread Hierarchy: Threads, Blocks, and Grids}

CUDA organises threads into three levels~\cite{nvidia2025cudaguide}: \textbf{threads} are the smallest execution units, each identified by \texttt{threadIdx}; \textbf{thread blocks} group threads that run on a single SM and share fast shared memory (up to 1024 threads per block)~\cite{nvidia2025cudaguide,sanders2010cuda}; \textbf{grids} consist of all blocks launched by a kernel, with ``each block of threads [...] scheduled on any of the available multiprocessors within a GPU, in any order''~\cite{nvidia2025cudaguide}. Each thread computes a global index:

\begin{equation}
\text{global\_idx} = \texttt{blockIdx.x} \times \texttt{blockDim.x} + \texttt{threadIdx.x}
\end{equation}

This maps each thread to a unique individual in the population array~\cite{nvidia2025cudaguide}.

\subsubsection*{Warps and SIMT Execution}

The hardware schedules threads in groups of 32 called warps - the basic scheduling unit on NVIDIA GPUs~\cite{nvidia2025cudaguide}. Block sizes should be multiples of 32, since ``when the total is not a multiple of 32, the last warp [...] will have some lanes that are unused throughout execution''~\cite{nvidia2025cudaguide}. When threads in a warp take different branch paths, the warp executes both paths sequentially with divergent threads masked - called warp divergence~\cite{sanders2010cuda,hwu2022pmpp}.

\subsubsection*{Memory Hierarchy}

\begin{figure}[H]
    \centering
    \includegraphics[width=0.85\textwidth]{diagrams/sec2/cuda_memory_hierarchy.png}
    \caption{CUDA memory hierarchy: scope and access latency by level. Adapted from~\cite{nvidia2025cudaguide}.}
    \label{fig:cuda_memory_hierarchy}
\end{figure}

CUDA provides several memory types~\cite{nvidia2024programming}: \textbf{global memory} has large capacity but high latency (~375--400 cycles)~\cite{jia2018volta}; \textbf{shared memory} is per-block with low latency (~19--28 cycles, up to 100 KB per SM on CC 8.6)~\cite{jia2018volta,nvidia2024programming}, suited for caching training data or distance matrices; \textbf{registers} are the fastest, private per-thread; \textbf{constant memory} is read-only and broadcast-optimized~\cite{nvidia2024programming}. The main bottleneck is global memory access, so coalesced access patterns - where warp threads access contiguous addresses - are important for throughput~\cite{nvidia2024cuda}. Section 3.1.1 describes the memory optimizations applied in this work.

\subsubsection*{Mapping Fitness Evaluation to CUDA}

Fitness evaluation maps naturally to CUDA because each individual's fitness is computed independently. One thread is assigned per individual: with population $P$ and block size $B$, the kernel launches $\lceil P/B \rceil$ blocks. Thread $i$ reads individual $i$, computes fitness, and writes the result. This reduces time complexity from $O(P)$ on CPU to $O(1)$ on GPU with sufficient parallelism.

\end{document}
